\section*{Capítulo \ref{ch:b3:groups} --- Teoría de grupos}

\begin{solution}{exo:b3:groups:1}
(a) Sea $[G:H] = 2$. Para $g \in H$, $gH = H = Hg$. Para $g \notin
H$: las dos clases laterales por la izquierda son $H$ y $gH$, luego
$gH = G \setminus H$; análogamente $Hg = G \setminus H$. Por tanto
$gH = Hg$ para todo $g$: $H
\trianglelefteq G$.

(b) Por (a), $G/H$ es un grupo de orden $2$; la clase $\bar x$ cumple
$\bar x^2 = \bar e$, es decir, $x^2 \in H$.

(c) Supongamos $H \leq A_4$ con $\abs H = 6$, luego de índice $2$. Por
(b), $\sigma^2 \in H$ para todo $\sigma \in A_4$. Todo $3$-ciclo es uno
de esos cuadrados: si $\sigma^3 = e$, entonces $\sigma =
\sigma^4 = (\sigma^2)^2$. Así que $H$ contiene los ocho $3$-ciclos de
$A_4$: $\abs H \geq 8 > 6$, contradicción. (El recíproco de Lagrange
falla a la primera ocasión: $6 \mid 12$.)
\end{solution}

\begin{solution}{exo:b3:groups:2}
Digamos $G/Z(G) = \langle gZ(G) \rangle$. Todo $x \in G$ se escribe
entonces $x = g^k z$ con $k \in \Z$, $z \in Z(G)$. Para $x = g^kz$, $y =
g^l z'$:
\[
xy = g^k z\, g^l z' = g^{k+l} z z' = g^l z'\, g^k z = yx,
\]
elementos centrales que conmutan con todo: $G$ es abeliano.

Orden $p^2$: $Z(G) \neq \{e\}$ (\cref{thm:b3:groups:pfixed}), luego
$\abs{Z(G)} \in \{p, p^2\}$. Si fuera $p$, entonces $G/Z(G)$ tendría
orden $p$ y sería cíclico, lo que forzaría a $G$ abeliano y a $Z(G) =
G$ de orden $p^2$ --- contradicción. Por tanto $Z(G) = G$.

No abeliano de orden $p^3$: el \emph{grupo de Heisenberg}
\[
H_p = \left\{ \begin{pmatrix} 1 & a & c\\ 0 & 1 & b\\ 0 & 0 & 1
\end{pmatrix} : a, b, c \in \mathbb F_p \right\}
\leq GL_3(\mathbb F_p),
\]
de orden $p^3$ (elección libre de $a, b, c$; el cierre y los inversos,
por cálculo directo). Es no abeliano: las dos matrices elementales
$I + E_{12}$ y $I + E_{23}$ tienen
\omterm{def:b3:groups:derived}{conmutador} $I + E_{13} \neq I$.
\end{solution}

\begin{solution}{exo:b3:groups:3}
(a) Un morfismo $f \colon \Z/n\Z \to \Z/n\Z$ queda determinado por $k =
f(\bar 1)$ (entonces $f(\bar m) = m\bar k$), y todo $\bar k$ define uno.
Es biyectivo si y solo si $\bar k$ genera $\Z/n\Z$, si y solo si
$\gcd(k,n) = 1$, si y solo si $\bar k \in (\Z/n\Z)^\times$. La
composición corresponde a la multiplicación: $f_k \circ f_l = f_{kl}$.
Por tanto
$\operatorname{Aut}(\Z/n\Z) \cong (\Z/n\Z)^\times$.

(b) La aplicación $\iota\colon G \to \operatorname{Aut}(G)$, $g \mapsto
\iota_g$, es un morfismo: $\iota_g \circ \iota_h = \iota_{gh}$. Su
imagen es $\operatorname{Inn}(G)$; su núcleo es $\{g : gxg^{-1} =
x\ \forall x\} = Z(G)$. El primer teorema de isomorfía da
$\operatorname{Inn}(G) \cong G/Z(G)$. Normalidad en
$\operatorname{Aut}(G)$: para $\alpha \in \operatorname{Aut}(G)$,
\[
(\alpha \circ \iota_g \circ \alpha^{-1})(x)
= \alpha\bigl(g\,\alpha^{-1}(x)\,g^{-1}\bigr)
= \alpha(g)\, x\, \alpha(g)^{-1} = \iota_{\alpha(g)}(x).
\]
\end{solution}

\begin{solution}{exo:b3:groups:4}
(a) Consideremos $\mu \colon H \times K \to HK$, $(h,k) \mapsto hk$,
sobreyectiva por definición. Fijemos $h_0k_0 \in HK$: entonces
$hk = h_0k_0
\iff h_0^{-1}h = k_0 k^{-1} \in H \cap K$. Escribiendo $u =
h_0^{-1}h$, la fibra de $h_0k_0$ es $\{(h_0u,\, u^{-1}k_0) : u
\in H \cap K\}$, de cardinal $\abs{H \cap K}$. Por tanto $\abs
H\,\abs K = \abs{H\times K} = \abs{HK}\,\abs{H \cap K}$.

(b) Si $HK$ es un subgrupo: $KH \subseteq HK$ porque $kh =
\bigl(h^{-1}k^{-1}\bigr)^{-1} \in (HK)^{-1} = HK$; y $HK
\subseteq KH$ tomando inversos en $HK = (HK)^{-1} \subseteq
(KH)^{-1}\dots$ --- más directamente, para $hk \in HK$, $(hk)^{-1} =
k^{-1}h^{-1} \in KH$, luego $HK = (HK)^{-1} \subseteq KH$; ambas
inclusiones dan $HK = KH$. Recíprocamente, si $HK = KH$: cierre,
$(hk)(h'k') = h(kh')k' \in h(HK)k' = (hH)(Kk') \subseteq HK$;
inversos, $(hk)^{-1} = k^{-1}h^{-1} \in KH = HK$; y $e \in HK$: es un
subgrupo. Si, por ejemplo, $K \trianglelefteq G$, entonces $hK = Kh$
para todo $h$, luego $HK = KH$ automáticamente.

(c) Para $h \in H$, $k \in K$, el
\omterm{def:b3:groups:derived}{conmutador} $[h,k] = hkh^{-1}
k^{-1}$ es igual a $(hkh^{-1})k^{-1} \in K$ ($K$
\omterm{def:b3:groups:normal}{normal}) y a $h(kh^{-1}k^{-1}) \in H$ ($H$
\omterm{def:b3:groups:normal}{normal}), luego está en $H \cap K =
\{e\}$: $hk = kh$.
\end{solution}

\begin{solution}{exo:b3:groups:5}
Contemos $E = \{(g,x) \in G \times X : g \cdot x = x\}$ de dos maneras:
\[
\abs E = \sum_{g \in G} \abs{\operatorname{Fix}(g)}
= \sum_{x \in X} \abs{G_x}
= \sum_{x \in X} \frac{\abs G}{\abs{\mathcal O_x}}
= \abs G \sum_{\mathcal O \text{ órbita}} \sum_{x \in \mathcal O}
\frac{1}{\abs{\mathcal O}}
= \abs G \cdot \#\{\text{órbitas}\},
\]
usando órbita--estabilizador ($\abs{G_x} = \abs G/\abs{\mathcal O_x}$) y
la partición en \omterm{def:b3:groups:action}{órbitas}.

Collares: sea $X$ el conjunto de las aplicaciones $\Z/p\Z \to
\{1, \dots, a\}$ (coloraciones de $p$ posiciones), $\abs X = a^p$, sobre
el que $\Z/p\Z$ actúa por rotación. La identidad fija las $a^p$
coloraciones. Una rotación $\bar k \neq \bar 0$ genera $\Z/p\Z$ ($p$
primo), de modo que una coloración fijada por ella es invariante por
\emph{todas} las rotaciones y, por tanto, constante: $a$ coloraciones
fijas. Burnside:
\[
\#\{\text{órbitas}\} = \frac{a^p + (p-1)a}{p} \in \N ,
\]
luego $p \mid a^p + (p-1)a$, es decir, $p \mid a^p - a$: el pequeño
teorema de Fermat, por puro recuento.
\end{solution}

\begin{solution}{exo:b3:groups:6}
Orden $15 = 3 \cdot 5$: $n_3 \mid 5$ y $n_3 \equiv 1 \pmod 3$ obligan a
$n_3 = 1$; $n_5 \mid 3$ y $n_5 \equiv 1 \pmod 5$ obligan a $n_5 = 1$.
Los \omterm{def:b3:groups:sylow}{subgrupos de Sylow} $P_3, P_5$ son
\omterm{def:b3:groups:normal}{normales}, se cortan trivialmente (órdenes
coprimos) y $\abs{P_3P_5} = 15$ (\cref{exo:b3:groups:4}(a)): por
la~\cref{prop:b3:groups:direct}, $G
\cong \Z/3\Z \times \Z/5\Z \cong \Z/15\Z$ (resto chino).

Orden $45 = 3^2 \cdot 5$: $n_3 \mid 5$, $n_3 \equiv 1 \pmod 3$ dan
$n_3 = 1$; $n_5 \mid 9$, $n_5 \equiv 1 \pmod 5$ dan $n_5 =
1$. Así que $G \cong P_3 \times P_5$ con $\abs{P_3} = 9 = 3^2$ y
$\abs{P_5} = 5$: ambos abelianos (el~\cref{thm:b3:groups:pfixed} para
$p^2$; el orden primo da cíclico), luego también lo es $G$.
\end{solution}

\begin{solution}{exo:b3:groups:7}
Orden $30 = 2 \cdot 3 \cdot 5$. $n_5 \mid 6$, $n_5 \equiv 1 \pmod
5$: $n_5 \in \{1, 6\}$; $n_3 \mid 10$, $n_3 \equiv 1 \pmod 3$:
$n_3 \in \{1, 10\}$. Supongamos $G$
\omterm{def:b3:groups:simple}{simple}, de modo que $n_5 = 6$ y $n_3 =
10$. Dos subgrupos distintos de orden primo $p$ se cortan trivialmente
(la intersección es un subgrupo propio de $\Z/p\Z$), así que los seis
$5$-subgrupos de Sylow aportan $6 \times 4 = 24$ elementos de orden $5$
y los diez $3$-subgrupos de Sylow aportan $10 \times 2 = 20$ elementos
de orden $3$: $24 + 20 = 44 > 30 - 1$ elementos distintos del neutro ---
absurdo. Luego $n_5 = 1$ o $n_3 = 1$: existe un \omterm{def:b3:groups:sylow}{subgrupo de Sylow}
\omterm{def:b3:groups:normal}{normal}.

Orden $56 = 2^3 \cdot 7$. $n_7 \mid 8$, $n_7 \equiv 1 \pmod 7$:
$n_7 \in \{1, 8\}$. Si $n_7 = 8$, los $7$-subgrupos de Sylow aportan $8
\times 6 = 48$ elementos de orden $7$, y quedan exactamente $56 - 48 =
8$ elementos más. Un $2$-subgrupo de Sylow tiene orden $8$ y consta de
esos elementos, luego es \emph{el} conjunto formado por ellos:
$n_2 = 1$. O bien $n_7 = 1$, o bien $n_2 = 1$: nunca es
\omterm{def:b3:groups:simple}{simple}.
\end{solution}

\begin{solution}{exo:b3:groups:8}
(a) La \omterm{def:b3:groups:action}{acción} $g \cdot xH = gxH$ da un
morfismo $\rho \colon G
\to \mathfrak S(G/H) \cong S_n$. Su núcleo es
\[
\ker\rho = \{g : \forall x \in G,\ gxH = xH\}
= \{g : \forall x,\ x^{-1}gx \in H\}
= \bigcap_{x \in G} xHx^{-1},
\]
un \omterm{def:b3:groups:normal}{subgrupo normal} (es un núcleo)
contenido en $H$ (tómese $x = e$). Si $N \trianglelefteq G$ y
$N \subseteq H$, entonces, para todo $x$: $N
= xNx^{-1} \subseteq xHx^{-1}$, luego $N \subseteq \ker\rho$: el núcleo
es el mayor de todos.

(b) Sea $[G:H] = p$ el menor primo que divide a $\abs G$, y sea $K =
\ker\rho \subseteq H$. Entonces $G/K$ se sumerge en $S_p$, luego $[G:K]$
divide a $p!$. Además $[G:K] = [G:H]\,[H:K] = p\,[H:K]$, así que $[H:K]$
divide a $(p-1)!$. Pero $[H:K]$ divide a $\abs G$, cuyos divisores
primos son todos $\geq p$, mientras que los divisores primos de $(p-1)!$
son todos $< p$: por tanto $[H:K] = 1$, es decir, $H = K = \ker \rho$ es
\omterm{def:b3:groups:normal}{normal}.

(c) Sea la \omterm{def:b3:groups:action}{acción} transitiva y sea
$x \in X$. La aplicación $\Phi
\colon G/G_x \to X$, $gG_x \mapsto g \cdot x$, está bien definida y es
biyectiva (órbita--estabilizador; la
\omterm{def:b3:groups:action}{órbita} es todo $X$), y entrelaza las
acciones:
$\Phi(h \cdot gG_x) = \Phi(hgG_x) = (hg)
\cdot x = h \cdot \Phi(gG_x)$.
\end{solution}

\begin{solution}{exo:b3:groups:9}
(a) $D(G)$ es \omterm{def:b3:groups:normal}{normal} con cociente
abeliano (la~\cref{prop:b3:groups:derived}); y si $N \trianglelefteq G$
tiene $G/N$ abeliano, la misma proposición da $D(G) \subseteq N$: $D(G)$
es el menor. Propiedad universal: sea $f \colon G \to A$ con $A$
abeliano. Entonces $f([x,y]) = [f(x), f(y)] = e$, luego $D(G)
\subseteq \ker f$, y el~\cref{thm:b3:groups:firstiso} factoriza $f =
\bar f \circ \pi$ a través de $G^{\mathrm{ab}}$, de manera única porque
$\pi$ es sobreyectiva.

(b) Los \omterm{def:b3:groups:derived}{conmutadores} son permutaciones
pares, luego $D(S_n) \subseteq A_n$. Recíprocamente, todo $3$-ciclo es
un \omterm{def:b3:groups:derived}{conmutador}:
\[
\bigl[(a\,b),\,(a\,c)\bigr] = (a\,b)(a\,c)(a\,b)(a\,c) =
(a\,b\,c),
\]
(comprobación directa en $a, b, c$), y los $3$-ciclos generan $A_n$
(el~\cref{lem:b3:groups:threecycles}): $D(S_n) = A_n$ para $n \geq 3$, y
$S_n^{\mathrm{ab}} \cong S_n/A_n \cong \Z/2\Z$. (Para $n = 2$: $S_2$ es
abeliano, $D(S_2) = \{e\}$, $S_2^{\mathrm{ab}} = S_2
\cong \Z/2\Z$ --- la fórmula $S_n^{\mathrm{ab}} \cong \Z/2\Z$ vale para
todo $n \geq 2$.)

$Q_8$: el cociente $Q_8/\{\pm 1\}$ tiene orden $4$, luego es abeliano y
$D(Q_8) \subseteq \{\pm 1\}$; y $[\mathrm i, \mathrm
j] = \mathrm i \mathrm j \mathrm i^{-1}\mathrm j^{-1} = \mathrm
i\mathrm j(-\mathrm i)(-\mathrm j) = (\mathrm i \mathrm j)^2 =
\mathrm k^2 = -1$, así que $D(Q_8) = \{\pm 1\}$ y
$Q_8^{\mathrm{ab}} \cong (\Z/2\Z)^2$ (orden $4$, exponente $2$: las
clases de $\mathrm i, \mathrm j$ tienen por cuadrado $\bar 1$).
\end{solution}

\begin{solution}{exo:b3:groups:10}
Inducción sobre $\abs G$; para $\abs G = p$ el único subgrupo propio es
$H = \{e\}$, y $N_G(\{e\}) = G \supsetneq
\{e\}$. Sea $Z = Z(G) \neq \{e\}$ (el~\cref{thm:b3:groups:pfixed}).

\emph{Si $Z \not\subseteq H$}: tómese $z \in Z \setminus H$; $z$ conmuta
con $H$, luego $zHz^{-1} = H$ y $z \in N_G(H) \setminus
H$.

\emph{Si $Z \subseteq H$}: pásese a $\bar G = G/Z$, un $p$-grupo de
orden menor, y $\bar H = H/Z \subsetneq \bar G$ (teorema de
correspondencia). Por inducción, $N_{\bar G}(\bar H) \supsetneq \bar H$;
tómese $\bar g \in N_{\bar G}(\bar H) \setminus \bar H$ y un
levantamiento $g$. Entonces $g \notin H$, y $gHg^{-1} \subseteq HZ = H$:
en efecto, $\overline{ghg^{-1}} = \bar g \bar h \bar g^{-1} \in \bar H$
significa $ghg^{-1} \in HZ = H$ (pues $Z \subseteq H$). Luego $g \in
N_G(H)\setminus H$.

Subgrupos maximales: si $M$ es maximal, $N_G(M) \supsetneq M$ obliga a
$N_G(M) = G$: $M \trianglelefteq G$. Entonces $G/M$ es un $p$-grupo sin
subgrupos propios no triviales (correspondencia y maximalidad). Tómese
$\bar x \neq \bar e$ en $G/M$, de orden $p^k$; entonces
$\bar x^{p^{k-1}}$ genera un subgrupo de orden $p$, que ha de ser todo
el grupo: $\abs{G/M} = p$.
\end{solution}

\begin{solution}{exo:b3:groups:11}
(a) $\abs{A_5} = 60$. Tipos de ciclo en $A_5$: $e$; dobles
trasposiciones, $\frac{1}{2}\binom{5}{1}\binom{4}{2}\cdot 1 = 15$ de
ellas ($5 \cdot 3$ maneras: elegir el punto fijo y después emparejar);
$3$-ciclos, $\frac{5 \cdot 4 \cdot 3}{3} = 20$; $5$-ciclos, $4! =
24$.

Una clase de $S_5$ contenida en $A_5$ o bien sigue siendo una sola
$A_5$-clase, o bien se desdobla en dos, según que el $S_5$-centralizador
de un elemento contenga o no una permutación impar
($\abs{\text{clase en }A_5} =
60/\abs{Z_{A_5}(\sigma)}$ y $Z_{A_5} = Z_{S_5} \cap A_5$). Para
$\sigma = (1\,2)(3\,4)$: $\abs{Z_{S_5}(\sigma)} = 120/15 = 8$, y
$(1\,2) \in Z_{S_5}(\sigma)$ es impar, luego $\abs{Z_{A_5}} = 4$ y la
clase tiene $60/4 = 15$ elementos: no hay desdoblamiento. Para $\sigma =
(1\,2\,3)$: $Z_{S_5}(\sigma) \supseteq \langle \sigma \rangle
\times \langle (4\,5)\rangle$, de orden $6 = 120/20$, luego son iguales;
contiene la permutación impar $(4\,5)$: clase de $60/3 = 20$: no hay
desdoblamiento. Para $\sigma$ un $5$-ciclo:
$Z_{S_5}(\sigma) = \langle \sigma\rangle$ (orden $120/24 = 5$), todas
pares: $Z_{A_5}(\sigma) = \langle
\sigma\rangle$ y la $A_5$-clase tiene $60/5 = 12$ elementos --- los $24$
ciclos de longitud cinco se desdoblan en \emph{dos} clases de $12$.
Cardinales de las clases: $1, 15, 20, 12, 12$.

(b) Un \omterm{def:b3:groups:normal}{subgrupo normal} $N$ es una unión
de \omterm{ex:b3:groups:actions}{clases de conjugación} que incluye
$\{e\}$, con $\abs N \mid 60$. Las sumas posibles $1 +
(\text{subconjunto de } \{15, 20, 12, 12\})$ son
\[
1,\ 13,\ 13,\ 16,\ 21,\ 25,\ 28,\ 28,\ 33,\ 36,\ 40,\ 40,\ 45,\
48,\ 48,\ 60 ;
\]
los únicos divisores de $60$ en la lista son $1$ y $60$: $N =
\{e\}$ o $A_5$.

(c) Sea $G$ \omterm{def:b3:groups:simple}{simple} con $\abs G = 60$.
$n_5 \mid 12$, $n_5
\equiv 1 \pmod 5$: $n_5 \in \{1, 6\}$. $n_5 = 1$ haría
\omterm{def:b3:groups:normal}{normal} el $5$-subgrupo de Sylow, en
contra de la simplicidad ($1 < 5 <
60$). Por tanto $n_5 = 6$.
\end{solution}

\begin{solution}{exo:b3:groups:12}
(a) $P$ es \omterm{def:b3:groups:normal}{normal} en $H = N_G(P)$ por
definición del \omterm{ex:b3:groups:actions}{normalizador}, y es un
$p$-subgrupo de Sylow de $H$ (su orden ya es la $p$-parte completa de
$\abs G$, con mayor razón de $\abs H$). Un
\omterm{def:b3:groups:sylow}{subgrupo de Sylow}
\omterm{def:b3:groups:normal}{normal} es único: cualquier otro sería
conjugado suyo (Sylow II en $H$) y, por tanto, igual a él.

(b) Sea $g \in N_G(H)$. Entonces $gPg^{-1} \subseteq gHg^{-1} = H$ es un
subgrupo de $H$ del mismo orden que $P$: un $p$-subgrupo de Sylow de
$H$, luego $gPg^{-1} = P$ por (a). Así $g \in
N_G(P) = H$: $N_G(H) \subseteq H$, y la inclusión recíproca es trivial.

(c) Supongamos $H \subseteq N \trianglelefteq G$ con $N$ propio. $P$ es
un $p$-subgrupo de Sylow de $N$; para cualquier $g \in G$,
$gPg^{-1} \subseteq N$ es otro, luego $gPg^{-1} = nPn^{-1}$ para algún
$n \in N$ (Sylow II en $N$), lo que da $n^{-1}g \in N_G(P)
\subseteq N$ y $g \in N$: $N = G$, contradicción (este es el
\emph{argumento de Frattini}). Para un subgrupo maximal $M \supseteq
N_G(P)$: $N_G(M) \supseteq M$ es $M$ o $G$; si $G$, entonces
$M \trianglelefteq G$ es un \omterm{def:b3:groups:normal}{subgrupo
normal} propio que contiene $N_G(P)$ --- excluido por el punto anterior.
Luego
$N_G(M) = M$.
\end{solution}

\begin{solution}{pb:b3:groups:1}
\textbf{1.} Para $x, y \in G$: $(xy)^2 = e$ da $xy = (xy)^{-1} =
y^{-1}x^{-1} = yx$ (cada elemento es su propio inverso): abeliano. Un
tal $G$, escrito aditivamente, es un espacio vectorial sobre $\mathbb
F_2$ ($2x = 0$, y los axiomas son los de grupo abeliano); si es finito,
tiene una base finita: $G \cong (\Z/2\Z)^k$, de orden
$2^k$.

\textbf{2.} Orden $p^2$: $G$ es abeliano
(el~\cref{thm:b3:groups:pfixed}). Si algún elemento tiene orden $p^2$,
$G$ es cíclico. En caso contrario, todos los $x \neq e$ tienen orden
$p$; aditivamente, $G$ es entonces un espacio vectorial sobre
$\mathbb F_p$ ($px = 0$), de dimensión $2$ ($p^2$ elementos):
$G \cong (\Z/p\Z)^2$.

Abelianos de orden $8$, según el orden máximo $m$ de un elemento: $m =
8$: cíclico $\Z/8\Z$. $m = 2$: $(\Z/2\Z)^3$ por la pregunta 1. $m = 4$:
sea $x$ de orden $4$ y $y \notin \langle x \rangle$; $y^2 \in
\langle x\rangle$ (índice $2$). $y^2 \in \{x, x^3\}$ daría a $y$ orden
$8$; luego $y^2 \in \{e, x^2\}$. Si $y^2 = x^2$, sustitúyase $y$ por
$xy$: $(xy)^2 = x^2y^2 = x^4 = e$ ($G$ abeliano) y $xy \notin
\langle x\rangle$. Podemos, pues, suponer $y^2 = e$: entonces $\langle x
\rangle \cap \langle y \rangle = \{e\}$, ambos
\omterm{def:b3:groups:normal}{normales} (abeliano),
$\abs{\langle x\rangle \langle y\rangle} = 8$:
la~\cref{prop:b3:groups:direct} da $G \cong \Z/4\Z \times \Z/2\Z$. Sin
repeticiones: los números de soluciones de $x^2 = e$ son $2, 4, 8$ en
los tres grupos.

\textbf{3.} Defínase $\psi \colon H \rtimes_{\varphi'} K \to H
\rtimes_{\varphi} K$ por $\psi(h, k) = (h, \alpha(k))$, que es una
biyección. Es morfismo:
\[
\psi\bigl((h,k)(h',k')\bigr)
= \bigl(h\,\varphi'(k)(h'),\, \alpha(kk')\bigr)
= \bigl(h\,\varphi(\alpha k)(h'),\, \alpha(k)\alpha(k')\bigr)
= \psi(h,k)\,\psi(h',k').
\]

\textbf{4.} $\operatorname{Aut}(\Z/n\Z) \cong (\Z/n\Z)^\times$
(\cref{exo:b3:groups:3}): para $n = 3$: $\{\pm 1\} \cong \Z/2\Z$;
$n = 4$: $\{\bar 1, \bar 3\} \cong \Z/2\Z$; $n = 5$: $\{\bar 1,
\bar 2, \bar 3, \bar 4\}$, cíclico de orden $4$ generado por $\bar
2$ ($2, 4, 3, 1$); $n = 7$: cíclico de orden $6$ generado por $\bar
3$ ($3, 2, 6, 4, 5, 1$). Para $V = (\Z/2\Z)^2$: un automorfismo es
$\mathbb F_2$-lineal (conserva la suma, y los escalares son $0,
1$), luego $\operatorname{Aut}(V) = GL_2(\mathbb F_2)$, de orden
$(4-1)(4-2) = 6$; actúa fielmente sobre los $3$ vectores no nulos, lo
que da un morfismo inyectivo en $S_3$ entre grupos de orden
$6$: $\operatorname{Aut}(V) \cong S_3$.

\textbf{5.} Cauchy proporciona $r$ de orden $p$; $N = \langle r
\rangle$ tiene índice $2$, luego es
\omterm{def:b3:groups:normal}{normal} (el~\cref{exo:b3:groups:1}).
Cauchy proporciona también $s$ de orden $2$, y $s \notin N$ (todos los
elementos de $N$ distintos del neutro tienen orden impar $p$).

\textbf{6.} $N \cap \langle s \rangle = \{e\}$ y $\abs{N\langle
s\rangle} = 2p$ (\cref{exo:b3:groups:4}(a)): por
la~\cref{prop:b3:groups:semidirect}, $G \cong \Z/p\Z \rtimes_\varphi
\Z/2\Z$ con $\varphi(1) = (x \mapsto sxs^{-1})$ un automorfismo de orden
divisor de $2$. En $(\Z/p\Z)^\times$, $k^2 = 1$ solo tiene las
soluciones $k = \pm 1$ ($X^2 - 1$ tiene a lo sumo dos raíces en el
cuerpo $\mathbb F_p$). Si $\varphi(1) = \mathrm{id}$: el producto es
directo, $G \cong \Z/p\Z \times \Z/2\Z \cong \Z/2p\Z$. Si
$\varphi(1) = -\mathrm{id}$: $G = \langle r, s \mid r^p = s^2 = e,
\ srs^{-1} = r^{-1}\rangle \cong D_p$
(\cref{ex:b3:groups:semidirectexamples}). No son isomorfos: $D_p$ es no
abeliano para $p \geq 3$ ($srs^{-1} = r^{-1} \neq r$).

\textbf{7.} $n_q \equiv 1 \pmod q$ divide a $p < q$: $n_q = 1$, luego
$N \cong \Z/q\Z$ es \omterm{def:b3:groups:normal}{normal}. Sea
$P \cong \Z/p\Z$ un $p$-subgrupo de Sylow: $N \cap P = \{e\}$, $NP = G$
(orden $pq$), luego $G
\cong \Z/q\Z \rtimes_\varphi \Z/p\Z$ con $\varphi \colon \Z/p\Z
\to \operatorname{Aut}(\Z/q\Z) \cong \Z/(q-1)\Z$ (cíclico, admitido). Si
$p \nmid q - 1$: la imagen de $\varphi$ tiene orden divisor de $p$ y de
$q-1$, luego es trivial, y $G \cong
\Z/pq\Z$ (\cref{ex:b3:groups:pq}). Si $p \mid q - 1$: aparte del trivial
$\varphi$, todo $\varphi$ no trivial es inyectivo (su núcleo, subgrupo
de $\Z/p\Z$, es trivial) y su imagen es \emph{el} único subgrupo $C$ de
orden $p$ del grupo cíclico $\Z/(q-1)\Z$. Dos
acciones no triviales $\varphi, \varphi'$
son entonces dos isomorfismos $\Z/p\Z \to C$, de modo que $\alpha =
\varphi^{-1}\circ\varphi' \in \operatorname{Aut}(\Z/p\Z)$ cumple
$\varphi' = \varphi\circ\alpha$: por la pregunta 3, los dos
\omterm{def:b3:groups:semidirect}{productos semidirectos} son isomorfos.
Hay, pues, exactamente un grupo no abeliano de orden $pq$ (no abeliano
porque $\varphi \neq
\mathrm{id}$ hace no trivial alguna conjugación). Orden $15$: $p =
3$, $q = 5$, $3 \nmid 4$: solo el cíclico.

\textbf{8.} No todo elemento tiene orden $\leq 2$ (si no, sería abeliano
por la pregunta 1) y ningún elemento tiene orden $8$ (si no, sería
cíclico y abeliano): algún $r$ tiene orden $4$, y
$N = \langle r\rangle$, de índice $2$, es
\omterm{def:b3:groups:normal}{normal}.

\textbf{9.} $s^2 \in N$ por el~\cref{exo:b3:groups:1}(b). El conjugado
$srs^{-1} \in N$ tiene orden $4$, luego $srs^{-1} \in \{r,
r^3\}$; si $srs^{-1} = r$, entonces $r$ y $s$ conmutan y $G =
\langle r, s\rangle$ sería abeliano --- excluido. Luego $srs^{-1} =
r^{-1}$. Si $s^2 = r$ o $r^3$, entonces $s$ tendría orden $8$: excluido.
(Alternativamente: $s^2$ conmuta con $s$, pero $s r s^{-1} =
r^{-1}$ y $s r^3 s^{-1} = r^{-3} = r$: ni $r$ ni $r^3$ quedan fijos al
conjugar por $s$.) Por tanto $s^2 \in \{e, r^2\}$.

\textbf{10.} Si $s^2 = e$: $G = \langle r, s \mid r^4 = s^2 = e,\
srs^{-1} = r^{-1}\rangle$. Los ocho elementos $r^is^j$ ($0 \leq i
< 4$, $0 \leq j < 2$) son distintos ($s \notin \langle r\rangle$) y las
relaciones determinan todos los productos: la asignación $r
\mapsto$ (rotación de $\pi/2$), $s \mapsto$ (una reflexión) define un
morfismo sobreyectivo sobre $D_4$, entre grupos de orden $8$: un
isomorfismo.

\textbf{11.} Si $s^2 = r^2$: de nuevo $G = \{r^i s^j\}$, y las
relaciones $r^4 = e$, $s^2 = r^2$, $srs^{-1} = r^{-1}$ fuerzan toda la
tabla. Con $\mathrm i = r$, $\mathrm j = s$, $\mathrm k =
rs$, $-1 = r^2$: $\mathrm i^2 = \mathrm j^2 = -1$, $\mathrm k^2 =
rsrs = r r^{-1} s s = s^2 = -1$ (usando $sr = r^{-1}s$) y
$\mathrm i \mathrm j \mathrm k = r\,s\,rs = r\,r^{-1}s\,s = s^2 =
-1$. Existencia: las matrices
\[
A = \begin{pmatrix} \iu & 0\\ 0 & -\iu \end{pmatrix},
\qquad
B = \begin{pmatrix} 0 & 1\\ -1 & 0 \end{pmatrix}
\]
cumplen $A^4 = I$, $B^2 = -I = A^2$ y $BAB^{-1} = A^{-1}$ --- para esta
última, compruébese
\[
BA = \begin{pmatrix} 0 & -\iu\\ -\iu & 0 \end{pmatrix} = A^{-1}B .
\]
Así pues, $\{\pm I, \pm A, \pm B, \pm AB\}$ es un grupo de orden $8$ que
realiza la tabla: $Q_8$ existe.

\textbf{12.} Sea $H \neq \{e\}$ un subgrupo y sea $x \in H
\setminus \{e\}$. Si $x \neq -1$, entonces $x \in \{\pm\mathrm i,
\pm\mathrm j, \pm\mathrm k\}$ y $x^2 = -1 \in H$. Así que $-1 \in H$
siempre. Los subgrupos son $\{e\}$, $\{\pm 1\}$ (el
\omterm{ex:b3:groups:actions}{centro}),
$\langle \mathrm i\rangle, \langle \mathrm j\rangle, \langle
\mathrm k\rangle$ (de índice $2$) y $Q_8$: todos
\omterm{def:b3:groups:normal}{normales} ($\{e\}$ y el
\omterm{ex:b3:groups:actions}{centro}, trivialmente; los de índice $2$,
por el~\cref{exo:b3:groups:1}; y $Q_8$ mismo). $D_4$ tiene cinco
elementos de orden $2$ ($r^2$ y las cuatro reflexiones) y $Q_8$ solo uno
($-1$): no son isomorfos. Un \omterm{def:b3:groups:semidirect}{producto
semidirecto} $H \rtimes K$ con $H, K \neq \{e\}$ exigiría $H \cap K =
\{e\}$, imposible porque ambos contienen $-1$.

\textbf{13.} $n_3 \mid 4$, $n_3 \equiv 1 \pmod 3$: $n_3 \in \{1,
4\}$; $n_2 \mid 3$, impar: $n_2 \in \{1, 3\}$. Si $n_3 = 4$: los cuatro
$3$-subgrupos de Sylow se cortan dos a dos trivialmente (orden primo),
lo que da $4 \times 2 = 8$ elementos de orden $3$; los $4$ elementos
restantes han de constituir el único $2$-subgrupo de Sylow: $n_2 = 1$.

\textbf{14.} $\ker\rho$ normaliza todo $3$-subgrupo de Sylow, luego
$\ker \rho \subseteq N_G(P_3)$, que tiene índice $n_3 = 4$, es decir,
orden $3$: $\abs{\ker\rho} \in \{1, 3\}$. El orden $3$ haría de
$\ker\rho$ un $3$-subgrupo de Sylow
\emph{\omterm{def:b3:groups:normal}{normal}}, en contra de $n_3
= 4$. Así que $\rho$ es inyectivo y su imagen $H \leq S_4$ tiene orden
$12$, índice $2$: $H \trianglelefteq S_4$ y $H$ contiene todos los
cuadrados (\cref{exo:b3:groups:1}(b)). Entre los cuadrados de $S_4$
están $e$ y los ocho $3$-ciclos ($\sigma = (\sigma^2)^2$ para un
$3$-ciclo), que generan $A_4$ (están en $A_4$ y, junto con sus
productos, dan los doce elementos; o bien:
el~\cref{lem:b3:groups:threecycles} para la parte de generación de
$n = 4$, que solo necesita $n \geq 3$). Luego $A_4 \subseteq H$ y
$\abs{A_4}
= \abs H$: $G \cong H = A_4$.

\textbf{15.} $P_3 \trianglelefteq G$, $P_3 \cap P_2 = \{e\}$,
$P_3P_2 = G$: $G \cong \Z/3\Z \rtimes_\varphi P_2$
(\cref{prop:b3:groups:semidirect}), $\varphi \colon P_2 \to
\operatorname{Aut}(\Z/3\Z) = \{\pm\mathrm{id}\} \cong \Z/2\Z$.
\begin{itemize}
  \item $\varphi$ trivial: los \omterm{prop:b3:groups:direct}{productos
  directos} $\Z/3\Z \times \Z/4\Z
        \cong \Z/12\Z$ y $\Z/3\Z \times (\Z/2\Z)^2 \cong \Z/6\Z
        \times \Z/2\Z$. \item $P_2 = \Z/4\Z$, $\varphi$ sobreyectivo:
  necesariamente $\varphi(1) = -\mathrm{id}$ (la única elección no
  trivial): un solo grupo, $\mathrm{Dic}_3 = \Z/3\Z \rtimes \Z/4\Z$.
  \item $P_2 = (\Z/2\Z)^2$, $\varphi$ sobreyectivo: $\ker\varphi$ es uno
  de los tres subgrupos de orden $2$; los tres $\varphi$ resultantes
  difieren en automorfismos de $(\Z/2\Z)^2$ que permutan esos subgrupos
  (pregunta 4: $\operatorname{Aut}
        \cong S_3$ actúa transitivamente sobre las tres involuciones),
  de modo que, por la pregunta 3, dan una única clase de isomorfía. Es
  $D_6$: tómese $t$ generador de $\ker\varphi$ y $x$ generador de
  $\Z/3\Z$; el elemento $\rho = (x, t)$ cumple $\rho^2 = (2x, 0)$,
  $\rho^3 = (0, t)$, $\rho^6 = e$ y ninguna potencia menor es $e$: orden
  $6$; y para $s = (0, u)$ con $u
        \notin \ker\varphi$: $s^2 = e$ y $s\rho s^{-1} = (-x, t)
        = \rho^{-1}$. Como $\langle \rho, s\rangle$ tiene orden $12$,
        $G \cong D_6$.
\end{itemize}

\textbf{16.} Contando los elementos de orden $2$: $\Z/12\Z$ tiene $1$;
$\Z/6\Z\times\Z/2\Z$ tiene $3$; $D_6$ tiene $7$ (seis reflexiones y el
giro de media vuelta $\rho^3$); $A_4$ tiene $3$; $\mathrm{Dic}_3$ tiene
$1$ (solo $(0, 2)$: un elemento $(h, k)$ con $k$ de orden $4$ en
$\Z/4\Z$ tiene orden $4$). Esto separa todos salvo las parejas
$\{\Z/12\Z, \mathrm{Dic}_3\}$ y $\{\Z/6\Z\times\Z/2\Z, A_4\}$: los
primeros miembros son abelianos y los segundos no ($\mathrm{Dic}_3$: la
\omterm{def:b3:groups:action}{acción} es no trivial; $A_4$: $(1\,2\,3)$
y $(1\,2)(3\,4)$ no conmutan). Cinco grupos distintos; las partes II--IV
muestran que la lista es completa.

\textbf{17.} La tabla de clasificación:
\begin{center}
\begin{tabular}{c l c}
$n$ & grupos de orden $n$ & \#\\
\hline
$1$ & $\{e\}$ & $1$\\
$2$ & $\Z/2\Z$ & $1$\\
$3$ & $\Z/3\Z$ & $1$\\
$4$ & $\Z/4\Z$, $(\Z/2\Z)^2$ & $2$\\
$5$ & $\Z/5\Z$ & $1$\\
$6$ & $\Z/6\Z$, $S_3 = D_3$ & $2$\\
$7$ & $\Z/7\Z$ & $1$\\
$8$ & $\Z/8\Z$, $\Z/4\Z{\times}\Z/2\Z$, $(\Z/2\Z)^3$, $D_4$,
      $Q_8$ & $5$\\
$9$ & $\Z/9\Z$, $(\Z/3\Z)^2$ & $2$\\
$10$ & $\Z/10\Z$, $D_5$ & $2$\\
$11$ & $\Z/11\Z$ & $1$\\
$12$ & $\Z/12\Z$, $\Z/6\Z{\times}\Z/2\Z$, $D_6$, $A_4$,
       $\mathrm{Dic}_3$ & $5$\\
$13$ & $\Z/13\Z$ & $1$\\
$14$ & $\Z/14\Z$, $D_7$ & $2$\\
$15$ & $\Z/15\Z$ & $1$\\
\end{tabular}
\end{center}
Los órdenes $6, 10, 14$ son la Parte II con $p = 3, 5, 7$; el orden $15$
es la pregunta 7; el orden $8$ es la Parte III junto con la pregunta 2;
el orden $12$ es la Parte IV; los órdenes primos son Lagrange; y los
órdenes $4$ y $9$ son la pregunta 2.

\textbf{18.} $Z = Z(G)$ es no trivial (el~\cref{thm:b3:groups:pfixed}) y
$Z \neq G$ (no abeliano), luego $\abs Z \in \{p, p^2\}$. Si
$\abs Z = p^2$, entonces $G/Z$ es cíclico de orden $p$ y
el~\cref{exo:b3:groups:2} hace $G$ abeliano: excluido, luego
$\abs Z = p$ y $\abs{G/Z} = p^2$. Por la pregunta 2, $G/Z$ es $\Z/p^2\Z$
o $(\Z/p\Z)^2$; el cíclico queda de nuevo excluido por
el~\cref{exo:b3:groups:2}: $G/Z \cong (\Z/p\Z)^2$. Como $G/Z$ es
abeliano, todo \omterm{def:b3:groups:derived}{conmutador} está en $Z$:
$D(G) \subseteq Z$; y $D(G) \neq \{e\}$ ($G$ no abeliano), luego
$D(G) = Z$ ($\abs Z = p$ no deja margen). Los
\omterm{def:b3:groups:derived}{conmutadores} son centrales y de orden
divisor de
$\abs Z = p$.

\textbf{19.} Inducción sobre $k$, siendo trivial el caso $k = 1$. Usando
$yx = xyz^{-1}\cdot$\,--- con precisión, $z = [y, x] =
yxy^{-1}x^{-1}$ da $yx = zxy$, es decir, pasar un $y$ hacia la izquierda
por delante de un $x$ produce un factor $z$, que es central y puede
aparcarse donde convenga. Entonces
\[
(xy)^{k+1} = (xy)^k\,xy = x^ky^kz^{k(k-1)/2}\,xy
= x^k\,(y^kx)\,y\,z^{k(k-1)/2}
= x^{k+1}y^{k+1}\,z^{k(k-1)/2 + k},
\]
puesto que llevar $x$ más allá de $y^k$ cuesta $k$ factores $z$ ($y^kx =
z^kxy^k$, por $k$ aplicaciones de $yx = zxy$); y $k(k-1)/2 + k
= k(k+1)/2$.

\textbf{20.} Con $k = p$: $(xy)^p = x^py^pz^{p(p-1)/2}$. Para $p$ impar,
$(p-1)/2$ es entero, luego $z^{p(p-1)/2} =
(z^p)^{(p-1)/2} = e$ (pregunta 18: $z$ tiene orden divisor de $p$):
$\theta(xy) = \theta(x)\theta(y)$, un morfismo. Sus valores son
centrales: la clase de $x$ en $G/Z \cong (\Z/p\Z)^2$ tiene orden divisor
de $p$, luego $x^p \in Z$. Si $\theta$ es trivial, todo elemento tiene
orden divisor de $p$: exponente $p$ (no $1$: $G \neq
\{e\}$). En caso contrario, hay algún $x^p \neq e$, y $x$ tiene orden
$p^2$ (su orden divide a $p^3$, y $x$ no puede tener orden $p^3$: $G$
sería cíclico y, por tanto, abeliano): exponente $p^2$.

\textbf{21.} Clases $\bar x, \bar y$ que generan $G/Z$: su
\omterm{def:b3:groups:derived}{conmutador} $z = [y, x]$ es $\neq e$,
pues en caso contrario $x, y, Z$ generarían un $G$ abeliano (sus clases
generan el cociente y $Z$ es central) --- y $z$ genera $Z$
($\abs Z = p$). Todo $g \in G$ tiene clase $\bar x^a\bar y^b$ para
$0 \leq a, b <
p$ únicos, luego $g = x^ay^bz^c$ con un $0 \leq c < p$ único: $p^3$
elementos, todos contabilizados. Los productos de estas formas
\omterm{def:b3:groups:normal}{normales} se calculan usando solo
$yx = zxy$, $z$ central y $x^p = y^p =
z^p = e$: la tabla queda forzada, de modo que dos grupos no abelianos
cualesquiera de orden $p^3$ y exponente $p$ son isomorfos (emparéjense
los generadores). El grupo de Heisenberg realiza las relaciones: con
$X = I + E_{12}$, $Y = I + E_{23}$, se calcula $[Y, X] = I -
E_{13}$ (central en $H_p$), y para cualquier $N$ estrictamente
triangular superior, $N^3 = 0$ da
\[
(I + N)^p = I + pN + \binom p2N^2 = I
\qquad\text{en característica } p,\ p \geq 3,
\]
puesto que $p \mid p$ y $p \mid \binom p2$ para $p$ impar: exponente
$p$. Así pues, el grupo de exponente $p$ es $H_p$.

\textbf{22.} $Z(G) = \langle r^p\rangle$: en efecto, $r^p$ es central
(el argumento de la pregunta 20: la clase de $r$ en el cociente $G/Z$,
de exponente $p$, da $r^p \in Z$) y es $\neq e$, luego genera el
\omterm{ex:b3:groups:actions}{centro}, de orden $p$. Tómese cualquier
$t \notin N$. Si $t^p = e$, póngase $s = t$. En caso contrario
$\theta(t) = t^p \in Z =
\langle r^p\rangle$, digamos $t^p = r^{pb}$; póngase $s = tr^{-b}$: por
la pregunta 20 ($\theta$ es un morfismo y $p$ es impar), $s^p =
t^pr^{-pb} = e$ y $s \notin N$. Conjugación: $srs^{-1} \in
N$ ($N$ \omterm{def:b3:groups:normal}{normal}) tiene orden $p^2$, luego
$srs^{-1} = r^m$ con $p
\nmid m$; además $s^p = e$ obliga a $m^p \equiv m \pmod{p^2}$ --- al
conjugar $p$ veces se recupera $r$, luego $m^p \equiv 1 \pmod
{p^2}$, y $m \equiv m^p \equiv 1 \pmod p$ (Fermat): $m = 1 +
ap$. La no trivialidad ($G$ no abeliano) da $a \not\equiv 0$;
sustituyendo $s$ por la potencia $s^{a'}$ con $aa' \equiv 1 \pmod p$, la
\omterm{def:b3:groups:action}{acción} se convierte en
$r \mapsto r^{1+p}$. Esto presenta $G$ como
$\Z/p^2\Z\rtimes_\varphi\Z/p\Z$ con $\varphi(1)\colon r
\mapsto r^{1+p}$; por la pregunta 3, dos morfismos no triviales
cualesquiera $\Z/p\Z \to \operatorname{Aut}(\Z/p^2\Z)$ con la misma
imagen --- y la imagen es \emph{el} único subgrupo de orden $p$ del
cíclico $\operatorname{Aut}(\Z/p^2\Z)$ --- dan
\omterm{def:b3:groups:semidirect}{productos semidirectos} isomorfos:
unicidad.

\textbf{23.} Abelianos: $\Z/p^3\Z$, $\Z/p^2\Z\times\Z/p\Z$, $(\Z/p\Z)^3$
(el argumento de la pregunta 2, un grado más arriba: clasificar según el
orden máximo). No abelianos: exactamente $H_p$ (exponente $p$, pregunta
21) y $\Z/p^2\Z\rtimes\Z/p\Z$ (exponente $p^2$, pregunta 22),
distinguidos por sus exponentes. Total: cinco. Para $p = 2$, el
argumento del morfismo de la pregunta 20 se viene abajo:
$z^{p(p-1)/2} = z^{1} = z \neq e$, elevar al cuadrado no es un morfismo
y, de hecho, los dos grupos no abelianos de orden $8$ tienen exponente
$4$ --- el invariante que separa $D_4$ de $Q_8$ es el número de
elementos de orden $2$ (cinco frente a uno), no el exponente. El mundo
de los $p$ impares es, por una vez, más ordenado que el de la
característica $2$.

\textbf{24.} En $H_p$ todo elemento $\neq e$ tiene orden $p$ (exponente
$p$, pregunta 21): $p^3 - 1$ elementos de orden $p$. En $M_p$, la
aplicación $\theta \colon x \mapsto x^p$ es un morfismo
$M_p \to Z(M_p) = \langle r^p\rangle$ (pregunta 20, $p$ impar);
$\theta(r) = r^p \neq e$, luego la imagen es todo el
\omterm{ex:b3:groups:actions}{centro}, de orden $p$, y
$\ker\theta = \{x : x^p = e\}$ tiene orden $p^3/p =
p^2$. Los elementos de orden $p$ son los elementos de ese núcleo
distintos del neutro: $p^2 - 1$. Para $p = 3$: $H_3$ tiene $27 - 1 = 26$
elementos de orden $3$ y $M_3 = \Z/9\Z \rtimes
\Z/3\Z$ tiene $9 - 1 = 8$. Para $p = 2$ el argumento muere de entrada:
elevar al cuadrado no es un morfismo en un grupo no abeliano de orden
$8$ (pregunta 23) y, en efecto, el conjunto $\{x : x^2 = e\}$ tiene $6$
elementos en $D_4$ --- que no es el orden de ningún subgrupo de $D_4$.
El recuento que sí separa la pareja es el número de elementos de orden
$2$: cinco en $D_4$ y uno en $Q_8$ (pregunta
11).

\textbf{25.} Si $p^2 \mid n$, el grupo $\Z/p\Z \times
\Z/(n/p)\Z$ tiene orden $n$ y no es cíclico: el orden de cada elemento
divide a $\operatorname{lcm}(p, n/p) = n/p < n$, pues $p \mid n/p$. Si
$p < q$ son primos que dividen a $n$ con $p \mid
q - 1$, la pregunta 7 proporciona un grupo no abeliano $\Z/q\Z \rtimes
\Z/p\Z$ de orden $pq$; entonces $(\Z/q\Z \rtimes \Z/p\Z) \times
\Z/(n/pq)\Z$ tiene orden $n$ y es no abeliano, luego no cíclico.
Supongamos ahora $\gcd(n, \varphi(n)) > 1$ y tomemos un primo $p$ que
divida a ambos. Escribiendo $\varphi(n) = \prod_{q^a \parallel
n} q^{a-1}(q - 1)$, la divisibilidad $p \mid \varphi(n)$ significa que,
o bien $p^2 \mid n$ (el factor $q^{a-1}$ con $q = p$, $a \geq 2$), o
bien $p \mid q - 1$ para algún primo $q \mid n$, $q
\neq p$: en ambos casos $n$ no es cíclico por lo anterior. Por
contraposición, $n$ cíclico obliga a $\gcd(n, \varphi(n)) = 1$.
Comprobación para $n \leq 15$: los valores $\varphi(n)$ para $n = 1,
\dots, 15$ son $1, 1, 2, 2, 4, 2, 6, 4, 6, 4, 10, 4, 12, 6,
8$, y $\gcd(n, \varphi(n)) = 1$ exactamente para $n = 1, 2, 3, 5,
7, 11, 13, 15$ --- exactamente las entradas de la tabla de la pregunta
17 con un solo grupo. Los demás órdenes se ven no cíclicos como antes:
$4, 8, 9, 12$ por un factor cuadrado y $6, 10, 12, 14$ por
$2 \mid q - 1$. (El recíproco --- $\gcd(n, \varphi(n)) = 1$ implica $n$
cíclico --- también es cierto; la pregunta 7 demuestra su primer caso no
trivial, $n = pq$ con
$p \nmid q - 1$.)
\end{solution}
